\documentclass[11pt]{article}
\usepackage[margin=1in]{geometry}
\usepackage{amsmath,amssymb,amsthm}
\usepackage[colorlinks=true,urlcolor=blue]{hyperref}
\newtheorem{theorem}{Theorem}
\newcommand{\vol}{\operatorname{vol}}
\title{Inherited-partition comparisons and a literature-subsumed target}
\author{Working note; no historical priority claim}
\date{November 15, 2026}
\begin{document}
\maketitle
Use one-sided conductance $\phi_G(S)=w(S,V\setminus S)/\vol_G(S)$.
The normalized Laplacian eigenvalues include $\lambda_1=0$.
The following theorem is a short composition of classical block
comparison and two-way Cheeger rounding. It is a sanity check for an
ambitious target below, not a new general Cheeger theorem.

\begin{theorem}[Pointwise leakage permits repair without a $k$ loss]
Let $G$ be a connected weighted loopless graph with positive degrees,
and let $C_1,\ldots,C_k$ be a partition of its vertices into nonempty
sets, $2\le k\le n$. Suppose that for every $v\in C_i$,
\[
 w_G(v,V\setminus C_i)\le h d_G(v).
\]
Let $G'$ result from inserting an edge of any weight $w>0$ between
two distinct vertices, or increasing an existing edge weight by $w$.
Then there is a partition of $V$ into exactly $k$ nonempty sets, each
of conductance at most
\[
 h+\sqrt{2\lambda_k(G')}.
\]
No gap between $\lambda_k$ and $\lambda_{k+1}$ is assumed.
\end{theorem}
\begin{proof}
If both endpoints lie in one old cluster, all old clusters still satisfy
the same pointwise leakage inequality, with the degrees of $G'$, so
the old partition has conductance at most $h$.
Otherwise merge the two endpoint clusters. This gives a partition
$S_1,\ldots,S_{k-1}$ whose pointwise external leakage remains at most
$h d_{G'}(v)$: the added edge is internal to the merged cluster,
external edge weights did not increase, and degrees did not decrease.

Put $D'=D_{G'}$, and let $L_{\rm int}$ contain just the Laplacians of
edges internal to these $k-1$ sets. Since the omitted cross-set edges
have positive semidefinite Laplacians,
\[
 (D')^{-1/2}L_{\rm int}(D')^{-1/2}
 \preceq (D')^{-1/2}L_{G'}(D')^{-1/2}.
\]
The left operator is block diagonal. Each block has a constant zero
mode (in its degree-weighted coordinate system), so its $k$th smallest
eigenvalue is the smallest second block eigenvalue; singleton blocks
have no second eigenvalue and can be assigned second eigenvalue
$+\infty$. At least one nonsingleton block exists because $n\ge k$.
Eigenvalue monotonicity now gives a block $S$ with degree-preserving
local second eigenvalue at most $\lambda_k(G')$.

Here ``degree-preserving local'' means the internal Laplacian on $S$
with the mass matrix $D'|_S$. It can be viewed as a weighted graph with
the external degrees replaced by holding weights. Two-way Cheeger
rounding, or its usual coarea proof with those masses, yields a
nontrivial split $S=A\sqcup B$ such that both pieces satisfy
\[
 \frac{w_{G'}(A,B)}{\vol_{G'}(A)},\quad
 \frac{w_{G'}(A,B)}{\vol_{G'}(B)}
 \le\sqrt{2\lambda_k(G')}.
\]
For clarity, choose the smaller-volume side in the scalar sweep;
the larger-volume side has the same cut numerator and at least as
large a denominator. If the local graph is disconnected, use a union
of components to obtain zero internal conductance instead.
The pointwise leakage inequality gives
$w_{G'}(A,V\setminus S)\le h\vol_{G'}(A)$ and the analogous bound
for $B$. Thus both original-graph conductances are at most
$h+\sqrt{2\lambda_k(G')}$. The remaining blocks have conductance
at most $h$. Splitting $S$ increases the number of nonempty parts
from $k-1$ to $k$.
\end{proof}

\paragraph{The average-leakage disjoint-set target is already subsumed.}
The initial candidate was to replace pointwise leakage by only
$\max_i\phi_G(C_i)\le h$ while asking for $k$ disjoint sets, without
requiring them to cover $V$, and retaining a bound
$C(h+\sqrt{\lambda_k(G')})$ independent of $k$.
The current primary source review found that this follows directly from
the core-refinement method in Lap Chi Lau's
\emph{Algorithmic Spectral Graph Theory}, draft September 16, 2026,
Section~7.7, pp.~99--104, Lemmas~7.18--7.21 and Algorithm~5.
The source develops that argument for regular graphs; replacing its
vertex counts by original degree volumes gives the weighted version
used here. No fractional-response argument is needed.
\url{https://cs.uwaterloo.ca/~lapchi/papers/ASGT.pdf}

Specifically, merge the two endpoint clusters to obtain $k-1$ sets
of conductance at most $h$ in $G'$. Set
\[
 \tau=\max\{h,4\sqrt{2\lambda_k(G')}\}.
\]
Trim each set to a nonempty minimum-conductance core, and add cores
from the remaining vertices while their conductance is at most $\tau$.
If this produces $k$ cores, the desired bound already holds. Otherwise
exactly $k-1$ cores remain. Lemma~7.20 retains at least $1/8$ of the
mass of every function in the low $k$-dimensional eigenspace on those
cores. Lemma~7.21 therefore gives at least $k$ local eigenvalues at most
$8\lambda_k(G')$ in their block diagonal local operator. One core has
local second eigenvalue at most $8\lambda_k(G')$; ordinary two-way
Cheeger gives two disjoint subsets of that core with local conductance
at most $4\sqrt{\lambda_k(G')}$. The core transfer inequality in
Lemma~7.18 gives global conductance at most
\[
 \boxed{\displaystyle
 \tau+8\sqrt{\lambda_k(G')}
 =\max\{h,4\sqrt{2\lambda_k(G')}\}
   +8\sqrt{\lambda_k(G')}.}
\]
Together with the remaining $k-2$ cores this gives exactly $k$
disjoint nonempty sets. For two endpoints in one old cluster, the old
partition itself has conductance at most $h$.
This is a specialization and degree-weighted transcription of a
published-on-the-web proof, not a new research frontier or breakthrough.
The full weighted reconstruction is retained in the sibling literature
directory in this session.

\paragraph{The remaining, narrower question.}
A potentially stronger question requires the new $k$ sets to cover
\emph{all} vertices, with maximum conductance
$C(h+\sqrt{\lambda_k(G')})$ and an explicit nearly-linear algorithm.
Core trimming leaves unassigned vertices, so the preceding corollary
does not itself supply this. Whether a known partition-completion theorem
settles that variant remains to be checked; it is an unverified candidate,
not a certified open problem. Algorithmic work must distinguish rounding
cost from constructing the fractional response factors.

\paragraph{Why it is a different target.}
The promise is the supplied old $k$-partition and a one-edge history.
In contrast, the general exact-$k$ theorem claimed by Y.~Li,
\emph{A sharp higher-order Cheeger inequality}, arXiv:2609.35369v2
(October 5, 2026), is
$\phi_k\le C\sqrt{\lambda_k\log(k+1)}$ with no such promise;
its noisy-hypercube example forces the logarithmic loss in general.
\url{https://arxiv.org/abs/2609.35369v2}.
That preprint was inspected as primary literature; no independent
full-proof certification is claimed here. The pointwise theorem above
explains one sufficient promise for a full partition, while the core
method already handles average leakage for disjoint sets. The separate
eight-vertex example in this directory proves that repair cannot in
general be confined to the two endpoint clusters.

\paragraph{Possible role for fractional response.}
Degree-aware relative response provides a small spectral certificate
for the edited operator and an exact logdet charge. It does not control
how a cut concentrates on old boundary vertices. A successful extension
would need a localization principle that turns that response certificate
into boundary control for the old clusters, or proves a dichotomy:
retain an old core or discover a sparse split elsewhere. Merely applying
Davis--Kahan to an old eigenspace would impose an additional eigengap
and would not settle this target.
\end{document}
